% GENERATED FILE. Edit canonical lessons, then run npm run generate:books.
% canonical-source-hash: fnv1a64:f501297bcf5fad29
% canonical-lessons: FR-C346-causer, FR-C346-provoquer, FR-C346-consequence, FR-C346-entrainer, FR-C346-decouler

\chapter{Causes et conséquences}
\label{ch:fr-causes-et-consequences}
\hlchaptermodality{346}

\begin{chapteropening}
\textbf{By the end of this chapter:} \emph{I can say what caused or provoked something, name a consequence, and say what it leads to or follows from.}

Complete the last lesson of chapter 346: I can say what caused or provoked something, name a consequence, and say what it leads to or follows from.
\end{chapteropening}

\section[causer]{causer — to cause}

\label{lesson:FR-C346-causer}



\begin{quote}
Before the new one: say the French for the opponent, then the French for to argue a case.
\end{quote}

\subsection*{You'll want to know: causer}
\textbf{causer} — \textquotedblleft{}to cause\textquotedblright{}.

\begin{quote}
Le brouillard a causé un accident sur l'autoroute.
\end{quote}

\begin{quote}
\emph{The fog caused an accident on the motorway.}
\end{quote}

Another \emph{causer}, informal, means to chat: \emph{causer avec un ami}.

\begin{cousinweb}
\emph{causer} is built on \emph{cause}, which comes from Latin \emph{causa} \textquotedblleft{}cause, reason, lawsuit\textquotedblright{} (EtymDB; Lewis \& Short).
\end{cousinweb}

\begin{grammarlens}[title={What you've built}]
One more word for telling, explaining and describing at B1.
\end{grammarlens}

\subsection*{Your turn}
\begin{itemize}
\raggedright
  \item \emph{Say it:} \emph{causer}
  \item \emph{Say it:} \emph{causer}, once more
  \item \emph{Say it:} a sentence of your own with \emph{causer}
  \item \emph{Recall:} say the French for the opponent, then the French for to argue a case, then say \emph{causer} again
\end{itemize}

\begin{tcolorbox}[breakable,colback=teal!4,colframe=teal!35!black,title={Before you move on}]
What does causer mean? (\textquotedblleft{}To cause\textquotedblright{}.) Say it once more.
\end{tcolorbox}
\section[provoquer]{provoquer — to bring about, to cause; to provoke}

\label{lesson:FR-C346-provoquer}



\begin{quote}
Before the new one: say the French for to argue a case, then the French for to cause.
\end{quote}

\subsection*{You'll want to know: provoquer}
\textbf{provoquer} — \textquotedblleft{}to bring about, to cause; to provoke\textquotedblright{}.

\begin{quote}
Ce médicament peut provoquer de la fatigue.
\end{quote}

\begin{quote}
\emph{This medicine can cause tiredness.}
\end{quote}

To cause, often something bad; also to provoke a person.

\begin{cousinweb}
\emph{provoquer} comes from Latin \emph{prōvocāre} \textquotedblleft{}to call forth, to challenge\textquotedblright{} (EtymDB; Lewis \& Short), built on \emph{vocāre} \textquotedblleft{}to call\textquotedblright{}.
\end{cousinweb}

\begin{grammarlens}[title={What you've built}]
One more word for telling, explaining and describing at B1.
\end{grammarlens}

\subsection*{Your turn}
\begin{itemize}
\raggedright
  \item \emph{Say it:} \emph{provoquer}
  \item \emph{Say it:} \emph{provoquer}, once more
  \item \emph{Say it:} a sentence of your own with \emph{provoquer}
  \item \emph{Recall:} say the French for to argue a case, then the French for to cause, then say \emph{provoquer} again
\end{itemize}

\begin{tcolorbox}[breakable,colback=teal!4,colframe=teal!35!black,title={Before you move on}]
What does provoquer mean? (\textquotedblleft{}To bring about, to cause; to provoke\textquotedblright{}.) Say it once more.
\end{tcolorbox}
\section[la conséquence]{la conséquence — the consequence}

\label{lesson:FR-C346-consequence}



\begin{quote}
Before the new one: say the French for to cause, then the French for to bring about.
\end{quote}

\subsection*{You'll want to know: la conséquence}
\textbf{la conséquence} — \textquotedblleft{}the consequence\textquotedblright{}.

\begin{quote}
Cette décision aura des conséquences pour tout le monde.
\end{quote}

\begin{quote}
\emph{This decision will have consequences for everyone.}
\end{quote}

\emph{En conséquence} — accordingly.

\begin{cousinweb}
\emph{conséquence} is a learned borrowing from Latin \emph{cōnsequentia} (EtymDB), from \emph{cōnsequī} \textquotedblleft{}to follow after\textquotedblright{}, built on \emph{sequī} \textquotedblleft{}to follow\textquotedblright{} (Lewis \& Short).
\end{cousinweb}

\begin{grammarlens}[title={What you've built}]
One more word for telling, explaining and describing at B1.
\end{grammarlens}

\subsection*{Your turn}
\begin{itemize}
\raggedright
  \item \emph{Say it:} \emph{la conséquence}
  \item \emph{Say it:} \emph{la conséquence}, once more
  \item \emph{Say it:} a sentence of your own with \emph{la conséquence}
  \item \emph{Recall:} say the French for to cause, then the French for to bring about, then say \emph{la conséquence} again
\end{itemize}

\begin{tcolorbox}[breakable,colback=teal!4,colframe=teal!35!black,title={Before you move on}]
What does la conséquence mean? (\textquotedblleft{}The consequence\textquotedblright{}.) Say it once more.
\end{tcolorbox}
\section[entraîner]{entraîner — to lead to, to bring about; to train}

\label{lesson:FR-C346-entrainer}



\begin{quote}
Before the new one: say the French for to bring about, then the French for the consequence.
\end{quote}

\subsection*{You'll want to know: entraîner}
\textbf{entraîner} — \textquotedblleft{}to lead to, to bring about; to train\textquotedblright{}.

\begin{quote}
Le froid entraîne souvent des rhumes.
\end{quote}

\begin{quote}
\emph{The cold often leads to colds.}
\end{quote}

In sport, \emph{s'entraîner} is to train; the coach is \emph{l'entraîneur}.

\begin{cousinweb}
\emph{entraîner} is built from \emph{en-} and \emph{traîner} \textquotedblleft{}to drag\textquotedblright{}, which is usually traced to a Vulgar Latin \emph{tragīnāre}, based on Latin \emph{trahere} \textquotedblleft{}to pull, to draw\textquotedblright{} (Lewis \& Short).
\end{cousinweb}

\begin{grammarlens}[title={What you've built}]
One more word for telling, explaining and describing at B1.
\end{grammarlens}

\subsection*{Your turn}
\begin{itemize}
\raggedright
  \item \emph{Say it:} \emph{entraîner}
  \item \emph{Say it:} \emph{entraîner}, once more
  \item \emph{Say it:} a sentence of your own with \emph{entraîner}
  \item \emph{Recall:} say the French for to bring about, then the French for the consequence, then say \emph{entraîner} again
\end{itemize}

\begin{tcolorbox}[breakable,colback=teal!4,colframe=teal!35!black,title={Before you move on}]
What does entraîner mean? (\textquotedblleft{}To lead to, to bring about; to train\textquotedblright{}.) Say it once more.
\end{tcolorbox}
\section[découler]{découler — to follow (from), to result from}

\label{lesson:FR-C346-decouler}



\begin{quote}
Before the new one: say the French for the consequence, then the French for to lead to.
\end{quote}

\subsection*{You'll want to know: découler}
\textbf{découler} — \textquotedblleft{}to follow (from), to result from\textquotedblright{}.

\begin{quote}
Tous nos problèmes découlent de cette erreur.
\end{quote}

\begin{quote}
\emph{All our problems follow from this mistake.}
\end{quote}

\emph{Découler de} — to flow from. Formal.

\begin{cousinweb}
\emph{découler} is built from \emph{dé-} \textquotedblleft{}from\textquotedblright{} and \emph{couler} \textquotedblleft{}to flow\textquotedblright{}, which is inherited from Latin \emph{cōlāre} \textquotedblleft{}to filter, to strain\textquotedblright{} (Lewis \& Short).
\end{cousinweb}

\begin{grammarlens}[title={What you've built}]
One more word for telling, explaining and describing at B1.
\end{grammarlens}

\subsection*{Your turn}
\begin{itemize}
\raggedright
  \item \emph{Say it:} \emph{découler}
  \item \emph{Say it:} \emph{découler}, once more
  \item \emph{Say it:} a sentence of your own with \emph{découler}
  \item \emph{Recall:} say the French for the consequence, then the French for to lead to, then say \emph{découler} again
\end{itemize}

\begin{tcolorbox}[breakable,colback=teal!4,colframe=teal!35!black,title={Before you move on}]
What does découler mean? (\textquotedblleft{}To follow (from), to result from\textquotedblright{}.) Say it once more.
\end{tcolorbox}
